% Portadilla y página de créditos de la edición impresa (compartida por los
% tres libros). La incluye build-print.sh con -B, después de definir:
%   \imprTitulo \imprSubtitulo \imprAutor \imprEditor \imprEdicion
%   \imprCita   \imprDOI       \imprISBN  \imprCopyright
% Un campo vacío se imprime como hueco marcado («PENDIENTE»): el piloto deja
% visible el lugar de lo que aún no existe en vez de inventarlo.
\newcommand{\imprHueco}[1]{%
  \fcolorbox{black!45}{black!6}{\parbox{0.86\linewidth}{\centering\footnotesize
    \textsc{pendiente}\enspace #1}}}
\newcommand{\imprCampo}[2]{\ifx#1\empty\imprHueco{#2}\else#1\fi}

% --- Portadilla (recto) ---
\begin{titlepage}
\centering
\vspace*{0.16\textheight}
{\Huge\bfseries\imprTitulo\par}
\vspace{0.8em}
{\Large\itshape\imprSubtitulo\par}
\vspace{3.2em}
{\large\imprAutor\par}
\vfill
% Lugar del DOI-concepto (Zenodo): uno por libro, estable entre versiones.
{\small\textsc{doi}\par}
\vspace{0.3em}
{\small\imprCampo{\imprDOI}{DOI-concepto del libro (Zenodo): 10.5281/zenodo.\_\_\_\_\_\_\_}\par}
\vspace{2.4em}
{\normalsize\scshape\imprEditor\par}
\end{titlepage}

% --- Créditos (verso) ---
\thispagestyle{empty}
\begingroup\small\raggedright
\vspace*{\fill}
\noindent\textbf{\imprTitulo}\\
\imprSubtitulo\\
\imprAutor

\vspace{1.2em}
\noindent\imprCampo{\imprCopyright}{Aviso de copyright: año y titular}

\vspace{0.8em}
\noindent\textbf{Edición:} \imprEdicion\\
\textbf{Editor:} \imprEditor\\
\textbf{DOI:} \imprCampo{\imprDOI}{DOI-concepto}\\
\textbf{ISBN:} \imprCampo{\imprISBN}{ISBN de esta edición impresa}

\vspace{1.2em}
\noindent Este libro se publica bajo la GNU Free Documentation License,
versión 1.3, publicada por la Free Software Foundation. Se permite copiarlo,
distribuirlo y modificarlo bajo los términos de esa licencia. Se incluye una
copia de la licencia en el apéndice titulado «GNU Free Documentation License».

\vspace{1.2em}
\noindent\textbf{Cómo citar:}\\
\imprCita

\vspace{1.2em}
\noindent agencydomains.org\\
github.com/gegolabs/agencydomains.org
\endgroup
\clearpage
